% TOP_PROBLEM: {"id": 30006743, "title": "Planted Clique Conjecture", "category_id": 15, "category": {"id": 15, "name": "computer_science", "display_name": "Computer Science", "description": "Computational complexity, algorithms, and theoretical CS.", "slug": "computer-science", "order_index": 15, "created_at": "2026-07-31T15:26:25.671Z"}, "status": "open"}
% FIRST_POSTED: 2026-09-27
% !TeX program = lualatex
% Compile twice: lualatex 30006743.tex
% All references and prose are embedded; no BibTeX database or external inputs.
\documentclass[11pt,a4paper]{article}
\usepackage[margin=24mm,headheight=14pt]{geometry}
\usepackage{fontspec}
\setmainfont{Latin Modern Roman}
\setsansfont{Latin Modern Sans}
\setmonofont{Latin Modern Mono}
\usepackage{amsmath,amssymb,mathtools}
\usepackage{unicode-math}
\setmathfont{Latin Modern Math}
\usepackage{microtype}
\usepackage{enumitem,xurl,needspace}
\usepackage[dvipsnames]{xcolor}
\usepackage{hyperref}
\hypersetup{unicode=true,colorlinks=true,linkcolor=MidnightBlue,urlcolor=MidnightBlue,
 pdftitle={500 Mathematical Problem Statements},pdfauthor={Source-annotated compilation},bookmarksnumbered=true}
\usepackage{bookmark}
\usepackage{tocloft}
\setlength{\cftsecnumwidth}{3em}
\cftsetpnumwidth{2.5em}
\cftsetrmarg{4em}
\usepackage{titlesec}
\titleformat{\section}{\Large\bfseries\raggedright}{\thesection}{0.7em}{}
\usepackage{fancyhdr}
\pagestyle{fancy}\fancyhf{}
\fancyhead[L]{\small\sffamily Mathematical problem statements}
\fancyhead[R]{\small Source edition 22}
\fancyfoot[C]{\thepage}
\setlength{\parindent}{0pt}
\setlength{\parskip}{5pt plus 1pt}
\setlength{\emergencystretch}{3em}
\setcounter{tocdepth}{1}
\setcounter{secnumdepth}{1}
\setlist[itemize]{leftmargin=3em,itemsep=5pt,topsep=4pt,parsep=0pt}
\allowdisplaybreaks
\newcommand{\N}{\mathbb{N}}
\newcommand{\Z}{\mathbb{Z}}
\newcommand{\Q}{\mathbb{Q}}
\newcommand{\R}{\mathbb{R}}
\newcommand{\C}{\mathbb{C}}
\newcommand{\F}{\mathbb{F}}
\newcommand{\A}{\mathbb{A}}
\newcommand{\PP}{\mathbb P}
\newcommand{\E}{\mathbb{E}}
\newcommand{\eps}{\varepsilon}
\newcommand{\dd}{\,\mathrm d}
\newcommand{\sref}[2]{\hyperlink{source:#1:#2}{[S#2]}}
\newcommand{\eref}[2]{\hyperlink{extra:#1:#2}{[E#2]}}
% Turn the next switch off for a shorter reading copy. The quotations preserve
% every exactTarget field, including qualifications omitted from a short summary.
\newif\ifshowcatalogue
\showcataloguetrue
\newcommand{\cataloguescope}[1]{\ifshowcatalogue
 \paragraph{Exact scope in the supplied catalogue.}
 \begin{quote}\small #1\end{quote}\fi}
\hypersetup{pdftitle={Problem 30006743 - Planted clique exact recovery},pdfauthor={Open Problems Math research notebook}}
\fancyhead[L]{\small\sffamily Open Problems Math \textbar\ Research notebook}
\fancyhead[R]{\small\sffamily 132 \textbar\ 27 September 2026}
\numberwithin{equation}{section}
\begin{document}
\setcounter{section}{0}
% ================================================================
% problemId: 30006743
% Canonical problem ID: 30006743
% exactTarget SHA-256: 0467ea804b5e8b27c747327862cbba100692a5f50a05149dd518cf07298e7db8
\clearpage
\section*{\texorpdfstring{Planted Clique Conjecture}{Planted Clique Conjecture}}\label{30006743}
\subsection{Definitions and mathematical statement}
Fix $0<\alpha<1/2$ and $k_n=\lfloor n^{1/2-\alpha}\rfloor$. Choose $S$ uniformly among $k_n$-subsets of $[n]$, place every edge inside $S$, and choose all other edges independently with probability $1/2$. The conjecture rules out a polynomial-time randomized estimator $A$ satisfying
\[
 \Pr[A(G)=S]\longrightarrow1.
\]
Probability includes the planted set, background edges, and algorithmic randomness. This is \emph{exact recovery}, not merely detection that a clique was planted. The algorithm may depend on $\alpha$ but has no length-dependent advice.
\par\smallskip\textit{Formulation and concept references:} \sref{30006743}{1}.
\cataloguescope{For every fixed real alpha with 0<alpha<1/2, set k\_n=floor(n\textasciicircum{}(1/2-alpha)). Sample a uniform k\_n-subset S of n labeled vertices, put every edge inside S, and independently include each other edge with probability1/2. There is no randomized algorithm running in time polynomial in n that, given this graph, outputs exactly S with probability tending to1 as n tends to infinity, over S, remaining edges and its internal coins. Algorithms may depend on the fixed alpha but receive no advice depending on n.}
\subsection{Short English statement}
Can one rule out efficient exact recovery of planted cliques whose size is a fixed power below the square-root threshold?
% BEGIN RESEARCH ADDITION 132
\subsection{Research attempt}

\paragraph{Current frontier and scope.}
The target is exact recovery of the hidden set for each fixed
$0<\alpha<1/2$, with no nonuniform advice.  The supplied paper studies
reductions among planted-clique conjectures, with parameters and success
notions made explicit; those reductions are not unconditional lower
bounds against all algorithms. \sref{30006743}{1}
As an adjacent restricted-model result, Mardia--Verchand--Wein analyze
nonadaptive low-degree detection in the \emph{above}-square-root,
sublinear-time regime.  That restriction and regime must not be
relabelled as a proof of the present below-square-root exact-recovery
claim. \eref{30006743}{1}  No unrestricted lower bound establishing the
catalogue target was verified in the literature check of
27 September 2026.

\paragraph{Attack map.}
A counterexample algorithm could try to find a small certified seed
inside the planted clique and then recover the rest by common neighbors.
Its two gaps are finding that seed efficiently and controlling adaptive
selection bias.  A hardness route can study the low-degree likelihood
ratio; the gap is transferring a fixed-degree calculation to every
uniform polynomial-time algorithm.  An information-theoretic route
might seek ambiguity between several large cliques, but the calculation
below rules that out in the selected regime.  All probabilities retain
the random planted set and independent background edges.

\paragraph{Attempt I: rule out an information-theoretic counterexample.}
Write $k=\lfloor n^{1/2-\alpha}\rfloor$ and condition on the planted
set $S$.  An alternative $k$-set $T$ with $|T\setminus S|=j$ is a
clique only if all
\[
 j(k-j)+\binom j2=jk-\frac{j(j+1)}2
\]
edges of $T$ not forced by the planting are present.  Their joint
probability is $2^{-jk+j(j+1)/2}$.  There are
$\binom kj\binom{n-k}j$ such sets.  Hence
\begin{align}\label{30006743:unique}
 \Pr[\exists\text{ another $k$-clique}\mid S]
 &\leq\sum_{j=1}^{k}\binom kj\binom{n-k}j
                  2^{-jk+j(j+1)/2}\\
 &\leq\sum_{j=1}^{k}\left(n^2 2^{-(k-1)/2}\right)^j=o(1).\notag
\end{align}
The last assertion holds because $k/\log n\to\infty$ for every
allowed fixed $\alpha$.  The bound is independent of the particular
$S$, so it remains valid without conditioning.  The hidden set is thus
the unique $k$-clique with high probability, and exhaustive search can
recover it.  That search has no polynomial runtime guarantee.
Ambiguity of the statistical model does not establish the conjectured
computational obstruction.

\paragraph{Attempt II: a correct seed-recovery lemma with fresh edges.}
Partition the vertices deterministically into two approximately equal
sets $U,W$.  Suppose a procedure inspects only $G[U]$ and outputs a
set $T\subset U$ of size
$t\geq(1+\varepsilon)\log_2 n$, for a fixed $\varepsilon>0$.
The crucial, still-unproved algorithmic requirement is that
$T\subseteq S$ with probability $1-o(1)$.

Conditional on $S$, the training information $G[U]$, and a correct
output $T$, the edges from $T$ to vertices of $W\setminus S$ are
fresh independent fair bits.  Define
\[
 C_W=\{v\in W:v\text{ is adjacent to every vertex of }T\}.
\]
All of $S\cap W$ belongs to $C_W$, while
\begin{equation}\label{30006743:fresh}
 \Pr[C_W\ne S\cap W,\ T\subseteq S]
 \leq n2^{-t}\leq n^{-\varepsilon}.
\end{equation}
This is a bound on the intersection event and does not assume the
correct-seed event has a probability bounded below.

With probability $1-o(1)$, hypergeometric concentration gives
$m=|S\cap W|\geq k/3$.  For each $v\notin S$, its number of neighbors
in the \emph{true} set $S\cap W$ is binomial with parameters $m,1/2$.
A Chernoff bound and a union bound show that, simultaneously for every
outside vertex, this number is less than $3m/4$ with failure at most
$n\exp(-c m)=o(1)$.  This is an unconditional high-probability event
conditional only on $S$; it need not stay independent of the event in
\eqref{30006743:fresh}.  A union bound is enough.

On both good events, output $C_W$ together with all vertices outside
$C_W$ adjacent to at least $3|C_W|/4$ of its vertices.  Every remaining
planted vertex is included and no outside vertex is included.  The
result is exactly $S$.  Therefore
\begin{equation}\label{30006743:seedrecovery}
 \Pr[\text{successful recovery}]
 \geq\Pr[T\subseteq S]-o(1).
\end{equation}
The cleanup takes polynomial time.  The entire unsolved algorithmic
burden of this route lies in finding a sufficiently large correct
seed from the training graph.

There is a real adaptivity pitfall here.  A claim that
\eqref{30006743:fresh} holds \emph{simultaneously for every}
$t$-subset $T\subseteq S$ selected after seeing all edges is false.
An outside vertex has approximately $k/2$ neighbors in $S$, much more
than $t$.  Choosing $T$ among those neighbors forces that vertex into
the common-neighbor set.  This is a counterexample to a universally
quantified seed lemma, not a proposed algorithm with access to the
hidden $S$.  The held-out set $W$ is what removes this bias in the
proved lemma.

For a uniformly random $t$-subset of $U$, the success probability is
\[
 \frac{(|S\cap U|)_t}{(|U|)_t}
 \leq\left(\frac{|S\cap U|}{|U|}\right)^t
 =\exp(-\Theta_\alpha((\log n)^2))
\]
with high probability, when $t$ is a fixed positive multiple of
$\log n$.  Here $(a)_t$ is the falling factorial.  Polynomially many
independent blind guesses therefore do not supply the missing seed.
Enumerating all such seeds is quasipolynomial rather than polynomial.

\paragraph{Attempt III: an exact fixed-degree likelihood calculation.}
Under the null model $P_0=G(n,1/2)$, let $X_e=1$ for a present edge
and $-1$ for an absent edge.  The likelihood ratio of the planted
model is exactly
\begin{equation}\label{30006743:likelihood}
 L(X)=\E_S\prod_{e\subset S}(1+X_e).
\end{equation}
For an edge set $F$, let $v(F)$ be the number of vertices incident to
its edges, and set $X_F=\prod_{e\in F}X_e$.  The monomials $X_F$ are
an orthonormal basis under $P_0$.  Expanding \eqref{30006743:likelihood}
gives the coefficient
\[
 \widehat L(F)=\Pr[V(F)\subseteq S]=\frac{(k)_{v(F)}}{(n)_{v(F)}}.
\]
For fixed integer $D\geq1$, let $L_{\leq D}$ be the projection onto
monomials using at most $D$ edges.  Then
\begin{equation}\label{30006743:lowdegree}
 \|L_{\leq D}-1\|_{L^2(P_0)}^2
 =\sum_{1\leq|F|\leq D}
       \left(\frac{(k)_{v(F)}}{(n)_{v(F)}}\right)^2
 \ll_D\sum_{v=2}^{2D}\left(\frac{k^2}{n}\right)^v=o(1),
\end{equation}
To justify the bound, for fixed $v\leq2D$ there are at most
$C_D n^v$ edge sets on their $v$ incident vertices with at most $D$
edges.  For large $n$, the falling-factorial ratio is at most
$C_D(k/n)^v$.  Finally $k^2/n\leq n^{-2\alpha}\to0$.

Thus every degree-at-most-$D$ real polynomial statistic $f$ normalized
by $\E_0f=0$ and $\E_0f^2=1$ satisfies
\[
 |\E_1f-\E_0f|
   =|\langle f,L_{\leq D}-1\rangle_{P_0}|
   \leq\|L_{\leq D}-1\|_{L^2(P_0)}=o(1).
\]
This is a rigorous mean-separation obstruction for fixed-degree
statistics.  It is \emph{not} a lower bound for all thresholdings of
such statistics, for degree growing arbitrarily with $n$, or for
all polynomial-time computations.  No equivalence between those
models has been proved in this notebook.

\paragraph{Stress test and dependency checks.}
Exact recovery would give a detector: check that the output is a
clique of size at least $\lceil(\log_2 n)^2\rceil$.  This check is
uniform and need not compute a possibly noncomputable fixed real
$\alpha$.  The planted output has that size eventually, whereas under
$P_0$ the probability of any such clique is at most
$\binom n\ell2^{-\binom\ell2}=o(1)$ for that $\ell$.
However, this implication does not convert the restricted
mean-separation estimate into a hardness theorem for every detector.
It is also not an assertion that every weak detection procedure
immediately gives exact recovery without additional reductions;
the distinctions made in \sref{30006743}{1} remain important.

The companion script verifies the Fourier coefficient and projected
norm identities by exact enumeration for a small graph.  It checks
the alternative-clique union-bound terms and the seed cleanup on
specified finite instances.  No finite sample success rate is
extrapolated to an asymptotic algorithm or lower bound.  The argument
allows constants and algorithms to depend on the fixed $\alpha$, but
adds no advice depending on $n$.

\paragraph{Outcome.}
\textbf{Outcome: Exploratory approach with unresolved gap.}

The planted set is shown to be information-theoretically identifiable;
a correctly trained logarithmic seed would yield efficient exact
recovery; and fixed-degree polynomial statistics have vanishing
normalized mean separation in the selected regime.  The adaptive-seed
shortcut is explicitly falsified.  Neither an efficient seed-finding
algorithm nor a bridge from the restricted likelihood calculation to
all uniform polynomial-time algorithms is supplied.  No resolution
or historical novelty claim is made.
% END RESEARCH ADDITION 132
\subsection{Sources}
\begin{itemize}\raggedright
\item[\textnormal{[S1]}]\hypertarget{source:30006743:1}{} Planted Clique Conjectures Are Equivalent, ECCC TR24-058,29March2024.
\url{https://eccc.weizmann.ac.il/report/2024/058/download/}.
{\small\textit{Catalogue formulation source.}}
% BEGIN RESEARCH REFERENCES 132
\item[\textnormal{[E1]}]\hypertarget{extra:30006743:1}{} Jay Mardia, Kabir Aladin Verchand, and Alexander S. Wein.
\emph{Low-degree phase transitions for detecting a planted clique in sublinear time}, Proceedings of the 37th Conference on Learning Theory, PMLR \textbf{247} (2024), 3798--3822. The model is nonadaptive low-degree detection in the above-square-root regime, not unrestricted exact recovery below that threshold.
\url{https://proceedings.mlr.press/v247/mardia24a.html}.
{\small\textit{Consulted 27 September 2026 for the indicated result or scope.}}
% END RESEARCH REFERENCES 132
\end{itemize}


\end{document}
